\chapter{The Trace}
\label{chap:the-trace}

The seam showed one pass becoming the next; but the machine, running pass after pass down the tower, does
more than hand each successor its residue. It keeps a record of what it did. That record is the trace, and
it is the object this chapter builds: a second tape, distinct from the one the passes are counted on, whose
entries are not passes but classes --- one cell for each rung of the descent, written on the way down from
the seam at the top to the bare difference at the bottom. The trace is built outright, cell by cell, and
certified in the building --- not a container whose contents are merely intended, but a record proved
together with everything written on it. It is the descent, written down.

\section{The second tape}
\label{se:the-second-tape}

The construction now carries two tapes, and they must never be confused. One is the record of revolutions:
a single cell per pass, the tally the seam increments, the count that will become a charge. The other is
the trace --- a single cell per class, the descent's own structure laid out rung by rung. The first counts
how many times the machine went around; the second records what each rung of one descent did. They are
different tapes with different alphabets, and the construction keeps their names apart precisely because a
reader who merged them would mistake the count of passes for the record of a pass, and lose both. This
chapter is the second tape: the trace, one entry per faculty, the machine's account of a single walk down.

\section{One cell per class}
\label{se:one-cell-per-class}

Each rung of the descent writes exactly one cell of the trace, and each cell holds a pair. On the way down,
a rung has two facts to record: what the pass built at that rung --- the constructed fact --- and what the
ground beneath it supplied --- the slip fact. The cell pairs them, the built against the given, so that the
trace at each rung is not a single reading but a comparison: what the machine made, set beside what it was
handed. This pairing is the whole content of a cell, and it is what makes the trace worth executing later:
a record of built-against-given, at every rung, is exactly what a test of whether the two agree will read.
The construction \emph{funges} the two facts of each rung into one cell under a shared tag, and it
\emph{tanges} the rungs apart by their place in the tower, so that the trace is an ordered stack of paired
cells, one for each class, and no cell is confused with its neighbour.

There is one rung that writes no cell. At the very top, at the seam, sits the anchor: it does not record a
comparison but seeds the tape the rest will write onto, holding the carried record so the first real cell
has something to strap itself to. So the trace has thirty-five cells --- one for each class from the true
output down to the distinguishing rung --- and, above them, the anchor that seeds them: thirty-six
declarations in all, of which thirty-five write and one carries. The distinction is not pedantry. The
anchor is where the trace of a later pass joins the trace it inherited; calling it a cell would double-count
the top of every descent.

\section{The bottom cell, and the tape that only appends}
\label{se:bottom-cell-append-only}

The trace is built downward, and it is built by appending only. Each rung straps its cell onto the record
the rung above it left --- never rewriting, never erasing, only adding one paired cell to the growing
stack. This is the same discipline the whole work has kept since the first mark: a record commits and does
not retract. The trace inherits it exactly, so that the finished stack is the honest history of the
descent, every rung's comparison preserved in the order it was made.

At the bottom sits the last cell, and it is the plainest of all. At the distinguishing rung --- the bare
difference where the tower began --- the constructed fact and the slip fact are each a single projection
deep: the origin fact the pass carried, set beside the origin fact the ground supplied. The bottom cell
traces the first difference directly, with nothing between the record and the fact it records. And because
the tape only appended on the way down, the finished trace is a stack that can be read back the other way
--- from the bottom up, the order a later chapter will need when it executes the trace rung by rung. This
is the tape the bootstrap's fifth chapter promised would one day be made to carry its own execution: it is
built here, and in the next chapter it is run.

\section{What is demonstrated, and what is assumed}
\label{se:demonstrated-assumed-III3}

What is demonstrated is the trace itself: a built record of one descent, thirty-five paired cells and the
anchor that seeds them, each cell pairing what the pass constructed against what the ground supplied, strung
into an append-only stack from the seam at the top to the origin fact at the bottom. Every cell is a
certified construction, and the trace as a whole is accepted, not intended --- the container and its
contents proved together, a record and not a promise.

What is assumed is unchanged and still small. The seed and its decidability; the honest import, marked as
borrowed; the single standing grant, the law of motion the first book named at its summit, still out of
play --- the trace records facts and pairs them, and reaches for no quantifier over infinity. What the
closing movement will weigh is the price of keeping this record along with everything else, and the answer
waiting there is short. With the descent written down --- one cell per class, built against given, the
first difference at the foot of the stack --- the construction can do the thing the trace was built for:
walk back up it and test, at each cell, whether the built and the given agree. That test is the next
chapter.
